\documentclass{muratcan_cv}

\setname{Michael}{Metternich}
\setaddress{Amsterdam/Netherlands}
\setmobile{+31 617128681}
\setmail{michael.j.metternich@gmail.com}
\setposition{Platform \& Infrastructure Engineering Leader} %ignored for now
\setlinkedinaccount{https://www.linkedin.com/in/michaelmetternich/}
\setgithubaccount{https://github.com/mjmetter}
\setthemecolor{blue}

% Deliberately no profile picture -- override the template's default header, which includes one.
\renewcommand\headerview{%
\begin{center}
    \name \\[0.05cm]
    \contact
\end{center}
}

\setcompanyname{Alpega Group}
\setcontactperson{Hiring Team}
\setclaimedposition{VP Software Engineering, Group}

\begin{document}
\headerview
\vspace{1ex}
\coverletter{ %
I'm writing to apply for the VP Software Engineering, Group role. Over eleven years at Bloomreach, I scaled platform, infrastructure, and, from 2025, application engineering, growing to roughly 45 engineers at peak, and I'm looking for the next scope. The mandate here, modernizing architecture, raising delivery velocity, and building a product-focused engineering culture, matches work I've already done at this altitude.

The platform-evolution and delivery-velocity goals map directly onto work I've already done. At Content, I built the first margin model for a platform with no per-customer visibility, consolidated two competing infrastructure generations onto a single EKS-based platform with full auto-scaling, and replaced fragmented logging with centralized observability; margin improved from under 50\% to 95\%+. At Discovery, I replaced reactive, on-request infrastructure work with a real platform: AWS account separation, Terraform-managed infrastructure, and a targeted monolith-to-Kubernetes migration assessed service by service rather than as a blanket rewrite. What particularly draws me to Alpega is the Group structure: TMS brings together products that started as separate companies onto one platform, and unifying previously separate engineering domains under a coherent technical foundation is exactly the kind of problem I love to work on.

Beyond infrastructure, at Content I also took over application engineering in early 2025 when the peer director who ran it departed, inheriting a team that had been overlooked for years under a struggling manager. I moved him onto a PIP, ultimately let him go, then hired and structured a new Director to fully own application engineering, reporting directly to me, unifying Content into a single organization spanning platform, infrastructure, and application development.

I've led AI adoption at this scale: I ran Bloomreach's company-wide AI coding-tool adoption, taking usage from an informal pilot to 160+ of 255 engineers within fifteen months, including deprecating legacy tooling and building the training curriculum that got teams there. I understand both the mechanics of driving that adoption and the cost side of it: what it takes to fund and govern AI tooling at scale, not just announce it.

On practicalities: the role is listed as open in Barcelona, Amsterdam, and Berlin. I'm Amsterdam-based and not looking to relocate to Spain, but I'm glad to travel frequently, which is close to how Discovery already worked across Amsterdam, Bangalore, and Mountain View. I'm currently on garden leave from August through November 2026 and available to start immediately after, though I'd welcome starting these conversations well before then.

I'd welcome the chance to talk through how my experience unifying platform and infrastructure across multiple business units and geographies maps onto what you're building across the Alpega Group.

\vspace{1ex}
Best regards, \\
Michael Metternich
}
\createfootnote
\end{document}
